\documentclass[12pt]{article}
\usepackage[utf8]{inputenc}
\usepackage[T1]{fontenc}
\usepackage{hyperref}
\title{Port Vato Language Resources for K-12 Curriculum Development}
\author{Compiled by Claude Sonnet 5.5 (claude-sonnet-5-5) for Andr\'as Kornai}
\date{}
\begin{document}
\maketitle

\section{Introduction}
Port Vato (ISO 639-3 ptv; Glottocode port1286), also called Daakie, is an Austronesian language of the Central Vanuatu branch, spoken on Ambrym Island, Vanuatu (about 1,300 speakers in 2001). The search results contained overview pages only; no school materials were found.

\section{Primary Educational Resources}
None found in the search results.

\section{Linguistic and Reference Materials}
\begin{itemize}
\item Wikipedia, ``Port Vato language'': \url{https://en.wikipedia.org/wiki/Port_Vato_language}
\item Wikipedia, ``Central Vanuatu languages'' (classification context): \url{https://en.wikipedia.org/wiki/Central_Vanuatu_languages}
\end{itemize}
The Wikipedia article mentions a DoReCo corpus (transcribed, time-aligned narrative recordings); I did not see a direct URL for it in the results.

\section{Religious and Community Resources}
\begin{itemize}
\item Joshua Project profile (a mission-oriented page; it says the group would benefit from a Bible in the language, and its population figures differ between versions): \url{https://joshuaproject.net/people_groups/14448/NH}
\end{itemize}

\section{Digital and Audio Resources}
Only the DoReCo corpus mentioned above; no direct link was found.

\section{Curriculum Development Context}
The language is described as stable but not known to be taught in schools. Little published material appears to exist, so a curriculum would depend on new documentation with the community and on the DoReCo recordings.
\end{document}